\section{Concluding Position}

LM-IDNet now contains a complete supplied-window detector implementation and a
substantial body of research evidence. The software processes device traffic,
fits and calibrates a DM model, exposes static and adaptive methods, validates
input compatibility, saves accepted state safely, restores and rolls it back,
and runs a bounded streaming service in a measured local Docker simulation.
The native likelihood work preserves the intended statistical objective and
records both numerical agreement and workload-dependent runtime.

The evidence also explains the system's limits. Protocol composition misses
some attacks. Legacy learning can improve one development sequence while
worsening another and admitting contaminated observations. Conservative guards
now prevent specified unsafe changes, but successful real-data adaptation is
not established by synthetic rejection or controlled promotion tests.
Temporal and multi-signal research provide further results without yet forming
one integrated reference implementation.

The project is ready for substantive thesis writing. Before finalising the
results, it still needs an explicit reference final-evaluation path and one
consistent, frozen evidence bundle, including the conservative comparisons.
The thesis should describe the delivered software, preserve useful negative
findings, and state exactly which claims belong to historical experiments,
current tests, or the agreed simulation. No additional hardware is needed to
write and evaluate that agreed scope.
